\chapter{Lymphoma Treatment}

\section*{Definition and Overview}
Lymphoma treatment refers to the comprehensive array of therapeutic interventions used to manage malignancies of the lymphatic system. These cancers are broadly divided into Hodgkin lymphoma (HL) and non-Hodgkin lymphoma (NHL). The primary objective of treatment is to achieve complete remission or cure by targeting and destroying clonal malignant lymphocytes, while minimising the risk of short-term and long-term toxicity. The selection of therapy is highly individualised, guided by the specific histological subtype, staging (commonly using the Lugano classification), molecular features, and patient-specific factors such as age, overall health, and cardiac or renal function.

\section*{Epidemiology}
Lymphomas represent the most common hematologic malignancies in Australia, collectively accounting for approximately 4\% of all new cancer diagnoses. Non-Hodgkin lymphoma is the more common form, typically diagnosed in older adults, with an incidence rate that increases with age and peaks in individuals in their 70s. In contrast, Hodgkin lymphoma exhibits a bimodal age distribution, most frequently affecting young adults aged 15 to 30, followed by a second peak in older adults aged 60 and above. Due to therapeutic advances, the 5-year survival rate for Hodgkin lymphoma in Australia is high at approximately 85\% to 90\%, while the 5-year survival rate for non-Hodgkin lymphoma stands at around 75\%.

\section*{Aetiology and Pathophysiology}
Lymphoma treatment modalities are directed at various biological mechanisms of the neoplastic lymphocytes to halt cell proliferation and induce apoptosis:
\begin{itemize}
    \item \textbf{Chemotherapy:} Systemic administration of cytotoxic drugs (e.g., the ABVD regimen for HL, or the CHOP regimen for NHL) that disrupt DNA replication and cell division in rapidly proliferating lymphoma cells.
    \item \textbf{Immunotherapy:} Monoclonal antibodies, such as rituximab, which targets the CD20 antigen on B-lymphocytes, triggering antibody-dependent cellular cytotoxicity and complement-mediated lysis.
    \item \textbf{Targeted therapies:} Small molecule inhibitors, such as Bruton's tyrosine kinase (BTK) inhibitors (e.g., ibrutinib, acalabrutinib) or B-cell lymphoma 2 (BCL-2) inhibitors (e.g., venetoclax), which block specific survival and signaling pathways in cancer cells.
    \item \textbf{Radiation therapy:} External beam radiotherapy directed at involved lymph node sites (involved site radiation therapy or ISRT) to cause DNA double-strand breaks, sterilising local disease.
    \item \textbf{Haemopoietic stem cell transplantation:} High-dose chemotherapy followed by autologous (patient's own) or allogeneic (donor) stem cell rescue, reserved for relapsed or refractory cases.
    \item \textbf{CAR T-cell therapy:} Chimeric antigen receptor T-cell therapy involves genetically modifying the patient's own T-cells to target specific B-cell antigens (e.g., CD19), offering a potential cure for refractory aggressive lymphomas.
\end{itemize}

\section*{Clinical Features}
The symptoms prompting treatment and monitored during therapy are reflective of lymph node enlargement, systemic cytokines, and bone marrow infiltration:
\begin{itemize}
    \item \textbf{B-symptoms:} Systemic symptoms including unexplained fever (temperature above 38 degrees Celsius), drenching night sweats, and significant weight loss (more than 10\% of body weight over 6 months).
    \item \textbf{Painless lymphadenopathy:} Firm, rubbery, mobile, and non-tender enlargement of lymph nodes in the cervical, axillary, or inguinal regions.
    \item \textbf{Fatigue:} Severe, persistent, and debilitating physical exhaustion that is not relieved by rest or sleep.
    \item \textbf{Pruritus:} Intense, generalised itching without a primary rash, particularly common in Hodgkin lymphoma.
    \item \textbf{Cytopenias:} Anaemia leading to dyspnoea and pallor, neutropenia leading to frequent infections, and thrombocytopenia causing petechiae or mucosal bleeding.
\end{itemize}

\section*{Differential Diagnosis}
Before initiating treatment, it is critical to confirm the diagnosis and rule out other causes of lymphadenopathy and systemic symptoms:
\begin{itemize}
    \item \textbf{Reactive lymphadenopathy:} Benign lymph node enlargement in response to viral (e.g., infectious mononucleosis, cytomegalovirus, HIV) or bacterial infections (e.g., tuberculosis, cat-scratch disease).
    \item \textbf{Other haematological malignancies:} Conditions such as acute or chronic leukaemias, or multiple myeloma, which may present with overlapping bone marrow and blood findings.
    \item \textbf{Metastatic solid tumours:} Carcinomas from primary sites like the breast, lung, colon, or head and neck that have spread to regional or distant lymph nodes.
    \item \textbf{Autoimmune and inflammatory disorders:} Conditions such as sarcoidosis, systemic lupus erythematosus (SLE), or rheumatoid arthritis, which can cause lymphadenopathy and systemic inflammatory features.
\end{itemize}

\section*{When to Seek Medical Care}
Patients undergoing lymphoma treatment are immunosuppressed and must monitor their symptoms closely.

\textbf{Call 000 (or local emergency services) for:}
\begin{itemize}
    \item \textbf{Neutropenic sepsis:} A life-threatening emergency in patients post-chemotherapy, characterised by a fever of 38 degrees Celsius or higher, rigours, chills, confusion, or rapid breathing.
    \item \textbf{Superior vena cava obstruction:} Swelling of the face, neck, or arms, combined with shortness of breath and blue discolouration of the skin, caused by mediastinal mass compression.
    \item \textbf{Severe bleeding:} Haemorrhage from the nose, mouth, or bowel that does not stop with direct pressure.
    \item \textbf{Acute shortness of breath:} Sudden chest pain or difficulty breathing, indicating pulmonary infection, bleomycin toxicity, or pulmonary embolism.
\end{itemize}

\section*{Diagnosis and Investigations}
A definitive diagnosis and pre-treatment staging are essential before formulating a management plan:
\begin{itemize}
    \item \textbf{Excisional lymph node biopsy:} The gold standard investigation; removing an entire intact lymph node allows histopathological analysis of tissue architecture and cellular morphology.
    \item \textbf{Immunohistochemistry (IHC) and flow cytometry:} Used to identify specific cell surface markers (such as CD20, CD30, CD15, and CD3) to differentiate and classify the lymphoma subtype.
    \item \textbf{Fluorodeoxyglucose (FDG) PET-CT scan:} The primary imaging modality for staging, treatment planning, and assessing response to therapy (using Deauville criteria).
    \item \textbf{Bone marrow aspirate and trephine biopsy:} Performed to evaluate bone marrow infiltration, which is particularly relevant in staging certain NHL subtypes.
    \item \textbf{Echocardiogram or MUGA scan:} Required prior to commencing anthracyclines (such as doxorubicin) to establish baseline left ventricular ejection fraction.
\end{itemize}

\section*{Management}
Management is multi-disciplinary and structured around chemotherapy, radiotherapy, and targeted agents, combined with supportive care:
\begin{itemize}
    \item \textbf{Induction therapy:} Combination chemotherapy regimens, such as ABVD (doxorubicin, bleomycin, vinblastine, dacarbazine) for HL, or R-CHOP (rituximab, cyclophosphamide, doxorubicin, vincristine, prednisolone) for diffuse large B-cell lymphoma.
    \item \textbf{Consolidation and radiotherapy:} Localised radiation therapy (ISRT) following induction chemotherapy for early-stage Hodgkin lymphoma or bulky disease.
    \item \textbf{Maintenance therapy:} Ongoing single-agent treatment (e.g., rituximab) in certain indolent lymphomas to prolong progression-free survival after successful induction.
    \item \textbf{Supportive medications:} Granulocyte colony-stimulating factors (G-CSF) to prevent neutropenia, antiemetics (e.g., ondansetron, aprepitant), and prophylactic antimicrobials (e.g., co-trimoxazole, aciclovir) to prevent opportunistic infections.
    \item \textbf{Survivorship and monitoring:} Regular clinical follow-ups, surveillance scans, and monitoring for late toxicities, including secondary cancers, cardiovascular diseases, and endocrine dysfunction.
\end{itemize}

\section*{Complications}
The treatments used for lymphoma carry significant potential risks and adverse effects:
\begin{itemize}
    \item \textbf{Tumour lysis syndrome (TLS):} A metabolic emergency caused by the rapid breakdown of cancer cells, leading to hyperkalaemia, hyperuricaemia, hyperphosphataemia, and acute renal failure.
    \item \textbf{Neutropenic sepsis and opportunistic infections:} High susceptibility to severe infections due to treatment-induced bone marrow suppression.
    \item \textbf{Organ toxicities:} Cardiotoxicity from anthracyclines (doxorubicin), pulmonary fibrosis from bleomycin, and peripheral neuropathy from vincristine.
    \item \textbf{Infertility:} Temporary or permanent gonadotoxicity, particularly associated with alkylating agents (e.g., cyclophosphamide), necessitating discussion of fertility preservation prior to treatment.
    \item \textbf{Second primary malignancies:} Increased risk of developing therapy-related myelodysplastic syndrome, acute myeloid leukaemia, or solid tumours in irradiated fields.
\end{itemize}

\section*{Prognosis and Outcomes}
The prognosis for patients undergoing lymphoma treatment is highly variable but generally positive. Hodgkin lymphoma is highly curable, with 5-year survival rates exceeding 90\% for early-stage disease and approximately 70\% to 80\% for advanced stages. Non-Hodgkin lymphoma prognosis depends heavily on the specific subtype: aggressive lymphomas (such as DLBCL) are curable in about 60\% to 70\% of patients with immunochemotherapy; indolent lymphomas (such as follicular lymphoma) are generally incurable but highly treatable, with patients often living for 10 to 20 years or more with sequential therapies.

\section*{Prevention and Self-Care}
While the occurrence of lymphoma cannot be directly prevented, patients can take self-care measures to manage their treatment and minimise complications:
\begin{itemize}
    \item \textbf{Infection control measures:} Avoiding crowds, washing hands frequently, wearing masks in high-risk areas, and avoiding raw or unpasteurised foods during periods of neutropenia.
    \item \textbf{Fever monitoring:} Keeping a digital thermometer at home and checking temperature immediately if feeling unwell, cold, or shivery.
    \item \textbf{Meticulous mouth care:} Using soft toothbrushes and alcohol-free mouthwashes to prevent and manage chemotherapy-induced oral mucositis.
    \item \textbf{Maintaining hydration:} Drinking plenty of water, particularly during and immediately after chemotherapy administration, to protect the kidneys.
\end{itemize}

\section*{Sources and Further Reading}
The following organisations provide evidence-based guidelines and support resources for patients undergoing lymphoma treatment:
\begin{itemize}
    \item Lymphoma Australia. \textit{Lymphoma Treatment.} https://www.lymphoma.org.au/lymphoma-treatment/
    \item Cancer Council Australia. \textit{Lymphoma.} https://www.cancer.org.au/cancer-information/types-of-cancer/lymphoma
    \item Healthdirect Australia. \textit{Lymphoma.} https://www.healthdirect.gov.au/lymphoma
    \item NHS. \textit{Non-Hodgkin lymphoma - Treatment.} https://www.nhs.uk/conditions/non-hodgkin-lymphoma/treatment/
    \item Leukemia \& Lymphoma Society (LLS). \textit{Treatment for Lymphoma.} https://www.lls.org/lymphoma/treatment
    \item National Cancer Institute (NCI). \textit{Lymphoma - Health Professional Version.} https://www.cancer.gov/types/lymphoma
\end{itemize}
